\documentclass{article}
\usepackage[utf8]{inputenc}
\usepackage{geometry}
\geometry{a4paper, margin=1in}
\usepackage{parskip}
\usepackage{hyperref}

\title{The Language of Collaboration: A Call for Precision}
\author{Albert Jan van Hoek}
\date{\today}

\begin{document}

\maketitle

\section*{The Language of Collaboration}

The most common words we use to call out bad collaboration are gendered slurs.

That’s not a coincidence. But it is a problem.

In my \textbf{Theory of Long-term Collaboration (TLC)}, collaboration is simply two learning networks trying to grow together. Anything irrelevant to that learning creates friction.

Gender is one of those irrelevancies — yet it’s baked into the very language we use when collaboration breaks down.

It’s time for a vocabulary upgrade.

Not to sanitize language. To make it more precise.

We need words that name what actually goes wrong — the behaviors that kill collaboration — without dragging in bodies, labels, or historical camps.

So on \textbf{International Women’s Day}, here’s a small act of linguistic resistance:

\begin{itemize}
    \item \textbf{Groovekiller} — destroys collaborative flow
    \item \textbf{Beatbreaker} — interrupts productive processes
    \item \textbf{Signaljammer} — blocks or distorts information
    \item \textbf{Vibehoarder} — takes credit without contributing
    \item \textbf{Floorfreezer} — makes others afraid to speak
    \item \textbf{Looplocker} — shuts down feedback and learning
    \item \textbf{Commonskipper} — avoids their share of maintaining the commons
\end{itemize}

\textit{``F-ing looplocker.''}

The behaviors are universal. The language should be too.

\section*{What would you add?}

\end{document}
